\subsection{Symmetry constraints and the exact scope of the LM4 comparisons}

Here \(F=\widehat F(x,y)\), \(z=(xy)^{-1}\), and
\(c(x)=x^{-1},c(y)=y^{-1}\). The semidirect product
\(E=F\rtimes\langle c\rangle\) is the free profinite product
of the three reflection groups represented by \(c,xc,y^{-1}c\).
We use a separate symbol \(\mathfrak B=C_2\amalg C_3\) for the
profinite modular group in the second half of this audit. No statement
below identifies an automorphism problem with a section problem.

\begin{result}{E.50}{Proved}{LM4 A.70: an exact finite source-central equivariance obstruction}
Let \(N=\ker(F\to(\mathbb Z/12)^2)\). The actual group
\(E/\overline{[N,E]}\), of order \(13824\), admits an involutory
automorphism fixing \(c\), preserving the three named reflection
classes, and preserving the three peripheral classes with common
multiplier congruent to seven modulo twenty-four. There is no full
reflection-class-preserving peripheral lift, of any order.
\end{result}
\begin{proof}
Write \(w=[x,y]=xyx^{-1}y^{-1}\). Since \(F'\subseteq N\),
centralizing \(N\) in \(F\) gives the class-two presentation with
\(w^{12}=1\). Its kernel over \((\mathbb Z/12)^2\) is generated
by \(x^{12},y^{12},w\). Centralizing it also with \(c\) imposes
\(x^{24}=y^{24}=1\), since \(c\) fixes \(w\). Conversely these
relations give the required central kernel, so the quotient has exactly
the normal forms \(x^ay^bw^dc^e\), with
\[
 a,b\in\mathbb Z/24,\quad d\in\mathbb Z/12,\quad e\in\mathbb Z/2,
\]
and multiplication
\[
 (a,b,d,e)(u,v,k,f)=
 (a+(-1)^eu,b+(-1)^ev,d+k-b(-1)^eu,e+f).
\]
This explicit group proves the converse as well as the order; a finite
count alone would not identify the actual source quotient.

The map \(\theta(x)=x^7,\theta(y)=y^7,\theta(c)=c\) is
\((a,b,d,e)\mapsto(7a,7b,d,e)\), because \(49\equiv1\pmod{12}\).
It squares to the identity and fixes the central kernel of order 48.
The conjugators for \(c,xc,y^{-1}c\) are \(1,x^3,y^{-3}\).
Moreover
\[
 (xy)^7=x^7y^7w^{-21},\qquad
 \theta(xy)=x^3(xy)^7x^{-3},\qquad
 \theta(z)=x^3z^7x^{-3}.
\]
The residue seven is a genuine profinite unit: take its dyadic component
to be \(-1\) and every odd component to be one.

Suppose a full lift existed. The parity character on the three
reflection classes preserves \(F\). Normalize the lift by an inner
automorphism from \(F\) so that it fixes \(c\); its common multiplier
\(\Lambda\) still satisfies \(\Lambda\equiv7\pmod{24}\).
If its restriction \(\phi\) sends \(x\) to
\(ax^\Lambda a^{-1}\), commutation with \(c\) gives
\(a^{-1}c(a)\in C_F(x)=\overline{\langle x\rangle}\).
The centralizer equality is the free-factor case of
\citet[Theorem B', pp.~159--160]{SymHerfortRibes1985}.
Writing \(a^{\rm ab}=(r,s)\) yields
\(a^{-1}c(a)=x^{-2r}\) and \(s=0\). Thus \(ax^{-r}\) is
fixed by \(c\). By \citet[Proposition 4(i), p.~574; proof p.~581]{SymLochakSchneps1997},
it is one. Therefore \(\phi(x)=x^\Lambda\), and similarly
\(\phi(y)=y^\Lambda\).

In \(S_4\), take \(X=(234)\), \(Y=(1234)\), with rightmost
permutations acting first. Since \(\Lambda\equiv7\pmod{12}\),
\(X^\Lambda Y^\Lambda\) has order two whereas
\((XY)^\Lambda\) has order four. This contradicts the third
peripheral condition. No order assumption on the full lift entered.
The independent certificate \path{code/check_symmetry_audit.py}
checks all 13824 normal forms, all 41472 generator edges, the central
kernel, the displayed conjugators and the two permutation orders.
\end{proof}

\begin{result}{E.51}{Proved}{LM4 A.74: prescribed frames can obstruct an otherwise trivial lift}
There is an actual finite quotient of \(E\) for which the identity
action, with multiplier one, has compatible prescribed frames on the
entire orientation-central refinement, including the required quotient
norms, but no full equivariant lift retaining those frames. Allowing
the frames to change restores the full identity lift.
\end{result}
\begin{proof}
In \(Q=\operatorname{SL}_2(\mathbb F_5)\), put
\[
 X=\begin{pmatrix}0&4\\1&4\end{pmatrix},\quad
 Y=\begin{pmatrix}0&1\\4&4\end{pmatrix},\quad
 J=\begin{pmatrix}0&1\\1&0\end{pmatrix}.
\]
Both \(X,Y\) have order three, generate \(Q\), and are inverted
by conjugation with \(J\). Generation also follows from
\(X^{-1}Y=-\left(\begin{smallmatrix}1&2\\0&1\end{smallmatrix}\right)\)
and \(XY^{-1}=-\left(\begin{smallmatrix}1&0\\2&1\end{smallmatrix}\right)\).
Their fifth powers give \(-I\), and upper and lower unipotents
generate \(Q\). Conjugating unipotents by \(\operatorname{diag}(2,3)\)
shows that \(Q\) is perfect. The independent matrix enumeration gives
orders120 for both the group and its derived subgroup.

For \(N=\ker(q:F\to Q)\), we claim
\[
 F/\overline{[N,F]}\xrightarrow{\ (q,\mathrm{ab})\ }
 Q\times\widehat{\mathbb Z}^{2}
 \quad\hbox{is an isomorphism}. \tag{E.51.1}
\]
The primary ATLAS entry records multiplier two for \(A_5\) and
identifies its double cover with \(\operatorname{SL}_2(5)\)
\citep{SymATLASA5}. As this double cover is perfect, it is the
universal central extension. Here is the needed centrally closed
consequence, including arbitrary profinite central kernels.
For a central extension \(D\to Q\) with kernel \(M\), perfectness
gives \(D=MD'\), hence the abstract derived subgroup \(D'\) is
perfect. An element of \(D'\) above the center of \(Q\) defines
a central-valued commutator homomorphism on \(D'\), which is zero.
Consequently \(D'\to A_5\) is central. Universality supplies a map
\(Q\to D'\) whose composite with \(D'\to Q\) is the identity.
This map is continuous because its domain is finite. Its image and the
central kernel generate \(D'\); perfectness makes the latter kernel
trivial. Thus \(D=Q\times M\), uniquely, and abelianization proves
(E.51.1). Under this identification
\[
 x=(X,e_1),\quad y=(Y,e_2),\quad
 c=(\operatorname{Ad}(J),-1).
\]
It follows that the actual full source-central quotient is
\((Q\times C_2^2)\rtimes C_2\), of order960.

Take the identity action and prescribe \(\beta=(-I,0)\) as the
\(x\)-cusp and \(xc\)-reflection conjugator, and one for the other
frames. This central, \(c\)-fixed element has order two, so all
conjugacy equations and \(\beta c(\beta)=1\) hold, even on the
entire group in (E.51.1).
For a hypothetical full lift, the external reflection frame has trivial
image in \(Q\). Normalize it to fix \(c\) literally. The normalizing
element lies in \(N\), so it does not alter any prescribed \(Q\)-frame
or the identity action in (E.51.1). Its restriction to \(F\) commutes
with \(c\), has multiplier one, and by the calculation in E.50 fixes
both \(x,y\). Every full conjugator of \(x\) then lies in
\(\overline{\langle x\rangle}\); its \(Q\)-image belongs to
\(\langle X\rangle\), of order three, and cannot be \(-I\).
This proves the prescribed-frame obstruction. It imposes the extra
equivariance equation, not merely a negative cocycle equation; the
identity with trivial frames remains an actual full solution.
\end{proof}

\begin{result}{E.52}{Proved}{LM4 B.43: a uniform forbidden simple section}
For every prime \(p\) and \(d\le97\), no subgroup of
\(\operatorname{GL}_d(\mathbb F_p)\) has a section isomorphic to
\(A_{10000}\). The same holds for every finite section of a closed
subgroup of \(\operatorname{Aut}(\widehat F_d^{(p)})\).
\end{result}
\begin{proof}
Pad the representation to dimension 97. If \(q\ne p\), a finite
\(q\)-subgroup lifts faithfully to \(\operatorname{GL}_{97}(\mathbb Z_p)\):
at each step the matrix lifting defect is an additive two-cocycle, and
averaging by the invertible group order kills its cohomology class.
The compatible corrections give the lift, faithful on reduction.
Over an algebraic closure of \(\mathbb Q_p\), its character is a sum
of ordinary irreducible characters, which also yields a complex
representation of dimension 97. Apply
\citet[Theorem A, p.~412]{SymCollins2007} in characteristic zero.
There is an abelian normal subgroup of index at most \(98!\).
Simultaneous diagonalization bounds the generator rank of every
subgroup and quotient of that abelian subgroup by 97.

If \(C_q^r\) is a section of the original \(q\)-group, intersect
with the abelian normal subgroup and then project. Its image has rank
at most 97, and the quotient has order at most \(98!\). Hence
\[
 r\le97+\lfloor\log_q(98!)\rfloor
 =608\quad(q=2),\qquad =419\quad(q=3).
\]
But \(A_{10000}\) contains \(C_2^{2500}\), from disjoint double
transpositions, and \(C_3^{3333}\), from disjoint three-cycles.
A Sylow subgroup of the inverse image of either subgroup under a
putative section map surjects onto it. Choose \(q=2\) unless \(p=2\),
and choose \(q=3\) in that case. Either rank is impossible.

The kernel of the Frattini action of the free pro-\(p\) automorphism
group is pro-\(p\). Indeed on successive characteristic finite
\(p\)-group quotients, the kernel is filtered by additive
\(\mathbb F_p\)-correction modules; take the inverse limit.
The finite groups having no \(A_{10000}\)-section are closed under
subgroups, quotients and extensions: in a putative simple section, the
image of the intersection with the normal subgroup is either trivial
or the whole simple group. Thus the section must already occur in one
of the two factors. This proves the automorphism assertion.
The two factorial logarithms are checked with exact integers in the
independent certificate, without a floating-point approximation.
\end{proof}

\begin{result}{E.53}{Proved}{LM4 B.42: an arithmetic ambiguity invisible on the Fermat relative tower}
Let \(F\) now be the index-six pure subgroup of
\(\mathfrak B=C_2\amalg C_3\), and put \(H_0=\overline{F'F^3}\).
The prosolvable residual \(R_{\rm sol}\) of \(G_{\mathbb Q}\) is
nontrivial. Its normalized arithmetic automorphisms, full cyclotomic
multiplier, and both pure torsion-generator conjugator frames are trivial
on \(\mathfrak B/R_{\rm nil}(H_0)\), including their modulo-twelve
images. Thus right multiplication by \(R_{\rm sol}\) preserves those
relative data for every full lift to which they are prescribed.
\end{result}
\begin{proof}
The cover for \(H_0\) is the degree-nine Fermat cubic
\(X^3+Y^3=Z^3\), with its nine coordinate-axis points removed.
Over \(k=\mathbb Q(\zeta_3)\) its boundary is rational, and the
automorphism \((X:Y:Z)\mapsto(\zeta_3X:Y:Z)\) fixes the origin
\((0:1:1)\). Its action on the compact Tate module has minimal
polynomial \(T^2+T+1\). Consequently the centralizer on the rational
Tate module is a commutative quadratic algebra, also at 2 and 3. In
\[
 0\longrightarrow\mathbb Z_p(1)^8\longrightarrow
 H_1(E^\circ_{\bar k},\mathbb Z_p)
 \longrightarrow T_pE\longrightarrow0,
\]
the two diagonal images of \(G_k\) are therefore abelian, and the
kernel of their product action is additive. The full homology image
is prosolvable. Extending from \(k\) to \(\mathbb Q\) adds only
an order-two quotient. The tangential action is this geometric action:
the branches \(X/Z=t^{1/3},Y/Z=(1-t)^{1/3}\), with leading coefficient
one, give the comparison; changing the lift changes homology by no
inner action.

Put \(P_p=\mathfrak B/O^p(H_0)\) and \(U_p=H_0^{(p)}\).
Here \(U_p\) is free pro-\(p\) of rank10 and has index54.
The deck action on its homology is faithful: triviality on the
rank-eight boundary module fixes all nine punctures, and a nonidentity
automorphism of a complex elliptic curve has at most four fixed points.
For the latter assertion, write it as \(z\mapsto az+b\): a nontrivial
translation has none, and for \(a\ne1\) the possible elliptic units
have \(\deg(a-1)\le4\).
The centralizer \(C_{P_p}(U_p)\) is therefore trivial, since its
intersection with the centerless free pro-\(p\) subgroup is trivial.
Restriction embeds \(\operatorname{Aut}(P_p;U_p)\) into
\(\operatorname{Aut}(U_p)\). The pro-\(p\) Frattini kernel and the
homology calculation show that the arithmetic image in this
automorphism group is prosolvable. The group \(P_p\) itself is
prosolvable as a pro-\(p\)-by-solvable group.

Write \(\phi_g(s)=\beta_gs\beta_g^{-1}\) and
\(\phi_g(r)=A_grA_g^{-1}\) on \(G_k\). Their pure frames are
unique, since the full free-factor centralizers intersect \(F\)
trivially, by \citet[Theorem B']{SymHerfortRibes1985}.
Composition gives
\(\beta_{gh}=\phi_g(\beta_h)\beta_g\), and the same for \(A\).
Thus \(g\mapsto(\beta_g^{-1},A_g^{-1},\phi_g)\) is a continuous
homomorphism to a prosolvable semidirect product
\((P_p\times P_p)\rtimes\rho_p(G_k)\).
The prosolvable residual of \(G_k\) equals \(R_{\rm sol}\):
intersect a solvable-quotient kernel of this normal index-two subgroup
with its conjugate, then extend by the order-two quotient. Therefore
\(R_{\rm sol}\) kills all these framed homomorphisms and the abelian
cyclotomic character.

The relative pronilpotent quotient is the fibre product of the
\(P_p\)'s over \(\mathfrak B/H_0\): its kernel is
\(H_0^{\rm nil}=\prod_pH_0^{(p)}\), so any tuple with one common
head is obtained by choosing a representative of that head and then
its independent kernel coordinates. This proves the simultaneous
assertion. For the marked modulo-twelve quotient the exact images are
\[
 q_{12}(F)=V_4\times\operatorname{SL}_2(\mathbb F_3),\qquad
 q_{12}(H_0)=V_4\times Q_8.
\]
Choose the projective sign by the diagonal-one representative modulo 4;
the four representatives form \(V_4\). The derived subgroup of
\(\operatorname{SL}_2(\mathbb F_3)\) is \(Q_8\), its abelian
quotient has order 3, and cubing preserves \(V_4\). The image of
\(H_0\) is a 2-group, so this quotient factors through \(P_2\).
Independent matrix enumeration verifies all four orders 96, 32, 24, 8
and the ambient order 576.

Finally \(x^5-x-1\) is irreducible modulo 3, and has square-free
factor degrees2 and3 modulo 2; both checks were executed. Its Galois
group is transitive of prime degree and contains a transposition by
cubing the latter Frobenius. It is primitive; the graph formed by the
conjugate transpositions is connected and its edges generate \(S_5\).
The image of \(R_{\rm sol}\) in this quotient is exactly \(A_5\):
it contains the solvable residual of \(S_5\) and is contained in the
sign kernel. Hence \(R_{\rm sol}\ne1\). The faithful arithmetic
embedding, stated in \citet[Theorem 1, p.~573]{SymLochakSchneps1997},
ensures that distinct such Galois elements are distinct full actions.
\end{proof}

\begin{result}{E.54}{Proved}{LM4 B.44: a simultaneous nontrivial kernel for two relative systems}
Let \(N_{12}=\ker(\mathfrak B\to
\operatorname{SL}_2(\mathbb Z/12)/\{\pm I\})\).
There is a nontrivial closed subgroup \(Z\le G_{\mathbb Q}\) whose
multiplier, normalized automorphism and both canonical pure frames
are trivial at every relative pro-\(p\) level over \(H_0\) and over
\(N_{12}\), simultaneously. Consequently a full lift with these
prescriptions leaves at least the entire coset \(\Phi Z\)
undetermined.
\end{result}
\begin{proof}
The two subgroups have indices 54 and 576, and free ranks 10 and 97.
Their finite quotient maps inject both finite free factors, so the
kernel is torsion-free by the full profinite torsion theorem
\citep[Theorem A']{SymHerfortRibes1985}; the open-subgroup Kurosh
formula, or Euler characteristic \(-1/6\), gives the ranks.
The actual arithmetic action preserves \(H_0\) and \(N_{12}\).
For the latter, the modular representation has conjugacy covariance
over \(\widehat{\mathbb Z}\), hence after reduction modulo 12;
\citet[\S4, pp.~774--776]{SymMarin2008} proves the covariance and
identifies its arithmetic shear. Conjugacy covariance preserves the
kernel, which is all that is needed here.

Let \(\mathscr C\) be the finite groups without an
\(A_{10000}\)-section. For \(U=H_0,N_{12}\), the group
\(P_{U,p}=\mathfrak B/O^p(U)\) is pro-\(\mathscr C\), being
an extension of the free pro-\(p\) group \(U_p\) by a group of
order at most 576. Its automorphism group preserving \(U_p\) is
also pro-\(\mathscr C\). Indeed restriction to \(U_p\) has finite
kernel of order at most \(576^{576}\): the centralizer of \(U_p\)
injects into the finite quotient, and an automorphism trivial on
\(U_p\) is determined by its corrections on at most 576 coset
representatives, each correction lying in that centralizer.
This order is smaller than \(10000!/2\). The restriction image is
pro-\(\mathscr C\) by E.52, so extension closure proves the assertion.

The class of pro-\(\mathscr C\) groups is closed under products
and closed subgroups as well. A finite quotient of a closed subgroup
of a product factors through a subgroup of a finite product of finite
quotients: take a product-open normal neighborhood whose intersection
lies in its kernel. Apply the finite subgroup and extension closures.
On \(G_k\), take the product of both framed homomorphisms of E.53
over both \(U\)'s and every \(p\), and retain the cyclotomic
character. Its image is pro-\(\mathscr C\). Let \(Z\) be its kernel.
All prescribed vanishing properties hold by construction.

To prove nontriviality, use the fully specified polynomial
\[
 f(X)=X^{10000}+15X^{9999}+10X^2+30.
\]
It is Eisenstein at 5. Modulo 2 its coprime factors are
\(X^{9999}(X+1)\); the lifted first factor is Eisenstein of odd
degree 9999, so tame inertia gives a 9999-cycle. Modulo 3 its coprime
factors are \(X^2(X^{9998}+1)\), with the latter separable. The
lifted quadratic factor is Eisenstein and tame, while the other roots
are unramified. Its inertia gives a transposition. Transitivity and
the 9999-cycle imply double transitivity; the conjugate-transposition
graph again proves that the group is \(S_{10000}\). These exact
conditions also verify \citet[Corollary 1 to Theorem 3, p.~442]{SymOsada1987}.

The image on \(G_k\) is \(S_{10000}\) or \(A_{10000}\). If
\(Z\) had trivial image, this finite quotient would factor through
the pro-\(\mathscr C\) image just constructed, a contradiction.
Its nontrivial normal image therefore contains \(A_{10000}\), so
\(Z\ne1\). Every factorization, order inequality and rank bound
used here is checked in \path{code/check_symmetry_audit.py}.
This proof constructs a nontrivial arithmetic kernel. Its effect on
the additional tripod and resolvent maps is determined next.
\end{proof}

\begin{result}{E.55}{Proved}{LM4 B.45: each tested value varies injectively on the arithmetic coset}
Use the subgroup \(Z\) of E.54. On the multiplier-one symmetric
peripheral group \(\mathcal H\), each labelled three-cusp test of
the icosahedral cover and each of the three correctly typed quintic
resolvents is injective on every coset \(\Phi Z\).
The first two resolvents have target \(\widehat{GT}_0\); the third
has target \(\operatorname{Out}_{\partial}(C_3\amalg C_3)\).
At most one member of \(\Phi Z\) belongs to the geometric kernel
\(\mathcal K\). The relative data therefore leave these values
undetermined, even when a full ordinary lift has already been chosen.
\end{result}
\begin{proof}
We reconstruct the maps needed here, so this assertion does not import
the older unexamined LM7 lifting theorem. Let \(q_5:\mathfrak B\to A_5\)
be the modular quotient and \(K=\ker q_5\). The quotient is the
triangle quotient obtained by imposing \(\tau^5=1\), where
\(\tau=sr\) is the cusp generator. Thus peripheral automorphisms
preserve \(K\). An automorphism with multiplier one induces an inner
automorphism of \(A_5\): \(\operatorname{Aut}(A_5)=S_5\), and an
odd automorphism interchanges the two classes of five-cycles. Adjust
its representative to induce the identity on \(A_5\). This adjustment
is unique modulo \(K\), because \(Z(A_5)=1\), and so gives a
canonical outer restriction to \(K\). A cusp conjugator has image
in the centralizer of the five-cycle, namely its cyclic subgroup.
Consequently this restriction fixes each of the twelve cusp classes,
indexed by \(A_5/\langle q_5(\tau)\rangle\).

The corresponding cover is \(X(5)\) with its twelve cusps removed.
Its compactification has genus zero, and its cusps are rational over
\(\mathbb Q(\zeta_5)\). Explicitly they are \(0,\infty\) and
\(\zeta_5^i\varphi,\zeta_5^i\bar\varphi\), with
\(\varphi=(1+\sqrt5)/2\); see
\citet[\S1.3]{SymFisher2003}. Killing the other nine cusp normal
subgroups gives the full fundamental group of the three-punctured
sphere. This follows directly from its profinite presentation, with
the product of the twelve meridians equal to one. Restriction and
quotient therefore give continuous homomorphisms
\(\mathfrak T_t:\mathcal H\to\mathcal P\).
For an arithmetic element fixing \(\mathbb Q(\zeta_5)\), the
construction is the natural arithmetic action on that three-punctured
curve. A fractional linear coordinate change over this field identifies
it with the usual tripod. Its Galois action is faithful, by the Belyi
faithfulness recorded in \citet[\S0.2 and Main Theorem(a)]{SymHarbaterSchneps2000}.
Every element of \(Z\) has trivial full cyclotomic character, so
this applies to \(Z\).

For the resolvents, let \(T=A_4\le A_5\) stabilize a letter, and
let \(J=q_5^{-1}(T)\). Since \(N_{A_5}(T)=T\), outer restriction
to \(J\) is canonical after an inner adjustment. The degree-five
cover has model
\[
 j=t^3(t^2+5t+40),\qquad
 j-1728=(t-3)(t^2+4t+24)^2.
\]
The branch cycles on five letters have types \(3\,1^2\),
\(2^2\,1\), and \(5\). Thus its orbifold has cone orders
\(2,3,3\), at \(3,r_+,r_-\), and one cusp, where
\(r_\pm=(-5\pm3\sqrt{-15})/2\).
Its full profinite group is \(C_2\amalg C_3\amalg C_3\).
For a multiplier-one action, normalize its \(A_5\)-action to the
identity. The residual conjugator of the three-cycle then belongs to
its centralizer in \(A_5\); this centralizer fixes the other two
letters separately. Hence both order-three classes are preserved.
The unique order-two class and cusp are preserved as well.
Killing either order-three factor gives a homomorphism to the
peripheral outer group of \(C_2\amalg C_3\), identified with
\(\widehat{GT}_0\) by the symmetric peripheral comparison of
\citet[\S1.3, Propositions 3--4]{SymHarbaterSchneps2000}.
Killing the order-two factor gives the third stated target, with both
order-three classes still named. Only this multiplier-one domain is
needed for B.45; no larger domain is assumed silently.

Put \(k=\mathbb Q(\sqrt{-15})\). The affine coordinates
\(j'_\pm=1728(t-r_\pm)/(3-r_\pm)\) identify the first two
filled orbifolds with the standard \((2,3,\infty)\) orbifold over
\(k\). Functoriality of arithmetic fundamental groups therefore
makes both resolvents the identity on the chosen Galois embedding of
\(G_k\). For the third, the coordinate
\((t-r_+)/(r_--r_+)\) gives the orbifold with cone orders three
at zero and one and a cusp at infinity. Its arithmetic action is
faithful on \(G_{k(\zeta_3)}\), as follows directly from the cover
\[
 \mathbb P^1_z\setminus\mu_3\longrightarrow
 \mathbb P^1_t(3,3;\infty),\qquad
 t=\frac{z^3}{z^3-1}.
\]
This is a degree-three finite \emph{etale map to the root stack},
with kernel \(P\cong\widehat F_2\) in
\(C_3\amalg C_3\to C_3\), sending the cone generators to
\(1,-1\). The rational point \(z=2\) gives compatible based
arithmetic actions. If an arithmetic element is inner on
\(C_3\amalg C_3\), its outer restriction to \(P\) is a deck
class. A nontrivial deck class permutes all three cusps \(\mu_3\),
whereas an element of \(G_{k(\zeta_3)}\) fixes them. Its action on
\(P\) is therefore outer trivial, and tripod faithfulness makes the
Galois element one. Again \(Z\subset G_{k(\zeta_3)}\).

Each map \(R\) just constructed is a homomorphism and is injective
on \(Z\). Thus \(R(\Phi\sigma)=R(\Phi)R(\sigma)\) has distinct
values for distinct \(\sigma\in Z\). A pure geometric mapping
class becomes trivial after forgetting to three punctures; the same
holds for its completion by continuity. Consequently
\(\mathfrak T_t(\mathcal K)=1\), and the injectivity just proved
allows at most one element of \(\Phi Z\) in \(\mathcal K\).
This proves the historical coset conclusion, but neither asserts that
the complete set of allowed lifts is one coset nor asserts that the
coset meets \(\mathcal K\).
\end{proof}

\begin{result}{E.56}{Proved}{The relative principal-congruence tower is not cofinal}
For \(12\mid N\), put \(U_N=\ker(\mathfrak B\to
\operatorname{SL}_2(\mathbb Z/N)/\{\pm I\})\).
The intersection \(\bigcap_{12\mid N}R_{\rm nil}(U_N)\)
is nontrivial. In particular compatibility of involutions on this
entire tower, as asserted in LM4 A.31, would not imply an involution
of the full profinite group. The proof of this limitation is independent
of the unexamined arithmetic existence clause of A.31.
\end{result}
\begin{proof}
Every finite quotient of \(\mathfrak B/R_{\rm nil}(U_N)\)
is an extension of a congruence quotient by a finite nilpotent group.
The only possible nonabelian simple composition factors are
\(\operatorname{PSL}_2(\mathbb F_p)\), with \(p\) prime:
reduce \(\operatorname{SL}_2(\mathbb Z/p^a)\) modulo \(p\),
whose congruence kernel is a \(p\)-group, and then remove its scalar
center. The same restriction holds for any finite quotient of the
open normal pure subgroup, by intersecting the finitely many conjugates
of its kernel in the ambient group.

The free rank-two pure group surjects onto \(A_7\): a seven-cycle
and a consecutive three-cycle generate it. The independent finite
enumeration in \path{code/check_symmetry_audit.py} produces all
2520 elements. This simple group is not any
\(\operatorname{PSL}_2(\mathbb F_p)\). Its order contains no
prime greater than seven, so such a \(p\) would be among
\(2,3,5,7\), whose projective special-linear orders are
\(6,12,60,168\), respectively.

If the intersection of the relative kernels were trivial, their
directed family would be cofinal among open normal subgroups: for
an open normal \(V\), compactness applied to the complements of
\(V\) gives finitely many relative kernels whose intersection lies
in \(V\), and a common multiple of their moduli gives one such
kernel. The \(A_7\)-quotient would then factor through a finite
quotient of the relative tower, contradicting its composition factors.
More strongly, the intersection maps onto \(A_7\), since its image
is normal and cannot be trivial. This is a failure of cofinality, not
a failure of compactness for a fixed family of closed conditions.
\end{proof}

\begin{result}{E.57}{Proved}{The ordinary central lifting criterion used by LM4's R2}
Let \(\Lambda\) be either the free profinite tripod group, with
peripherals \(x,y,(xy)^{-1}\), or \(C_2\amalg C_3\), with
peripherals \(s,r,(sr)^{-1}\). Let \(N\) be closed normal.
A multiplier-one peripheral automorphism of \(\Lambda/N\), with
three specified conjugators, has a full peripheral lift retaining
those conjugators if and only if it has such a lift to
\(\Lambda/\overline{[N,\Lambda]}\).
In particular every such automorphism of
\(\Lambda/R_{\rm nil}(U)\), for an open normal \(U\), lifts.
This assertion imposes no order, equivariance, or norm condition on
the full lift. It independently reconstructs the finite count and
compactness steps of LM7\_PUSH D.50 and E.58 used as LM4's R2.
\end{result}
\begin{proof}
We prove the finite count, including all proper subgroup images.
Fix a finite source quotient \(q:\Lambda\twoheadrightarrow E\),
write \(A=q(N)\), \(D=E/A\), and let \(T_i\) be the
\(A\)-conjugacy class of the three peripheral images \(X_i\).
For every subgroup \(J\le E\) onto \(D\), put \(M=J\cap A\).
We need only consider those \(J\) for which all \(T_i\cap J\)
are nonempty. For a \(J\)-invariant irreducible complex representation
\(\rho\) of \(M\), of dimension \(d\), put \(e=\exp(J)\) and
\[
 \widetilde J_\rho=\{(g,B):g\in J,\quad
 B\rho(m)B^{-1}=\rho(gmg^{-1}),\quad\det B\in\mu_e\}.
\]
Schur's lemma makes each fibre a torsor under \(\mu_{de}\).
Hence this is a finite central extension of \(J\). The normal graph
\(m\mapsto(m,\rho(m))\) meets the scalar kernel trivially.
Quotienting by it gives a central extension \(Q_\rho\to D\),
and the displayed group identifies exactly with
\(J\times_DQ_\rho\). Its matrix coordinate defines a representation
\(\Psi_\rho\) of this pullback restricting to \(\rho\) on \(M\).

There is an actual map \(R_\rho^0:\Lambda\to Q_\rho\) above
the given map to \(D\). For a free source choose two lifts. For the
modular source, choose \(a_s\in T_1\cap J\) and
\(a_r\in T_2\cap J\). Their implementing matrices have scalar
squares and cubes; scale them to have square and cube one. Their
determinants then lie in \(\mu_e\), so their classes give a map
from \(C_2\amalg C_3\), retaining its actual torsion relations.
Take the joint image of the maps to \(D\) and all \(Q_\rho\).
This gives an actual finite source quotient \(D^\dagger\to D\)
with central kernel. In particular its source map kills
\([N,\Lambda]\).

Here is the count which this refinement controls. For a surjection
\(\chi:\Lambda\to D\), and \(M\)-conjugacy-invariant subsets
\(U_i\subset J_{\chi(x_i)}\), use any actual lifts
\(R_\rho:\Lambda\to Q_\rho\) of \(\chi\). The number of
peripheral homomorphisms to \(J\) with these boundary values is
\[
 \frac1{|M|}\sum_{\rho\in\operatorname{Irr}(M)^J}
 \frac1{\dim\rho}
 \prod_{i=1}^3\sum_{a\in U_i}
             \operatorname{tr}\Psi_\rho(a,R_\rho(x_i)).
 \tag{E.57.1}
\]
For a direct proof, put \(B_i=\sum_{a\in U_i}a\in\mathbb C[J]\).
Each \(B_i\) commutes with \(\mathbb C[M]\), and
\(B_1B_2B_3\in\mathbb C[M]\). If one quotient generator moves
the primitive central idempotent of \(\rho\), the corresponding
block of \(B_i\) is zero. The first two quotient values generate
\(D\), so only \(J\)-invariant representations survive.
For one of these, each lifted matrix sum commutes with \(\rho(M)\)
and is its trace divided by \(d\), times the identity. The product
of three lifted elements has \(Q_\rho\)-coordinate one, and is
therefore the graph of their product in \(M\). The regular-character
identity \(\mathbf1_{\{1\}}(m)=|M|^{-1}\sum_\rho d\,
\operatorname{tr}\rho(m)\) now gives (E.57.1): three scalar
denominators, one matrix trace, and the factor \(d\) leave \(d^{-1}\).
This is a genus-zero argument. No positive-genus counting theorem is
being applied to the tripod.

Suppose a second epimorphism to \(D^\dagger\) conjugates its three
source peripherals by \(r_i\). Write \(h_i\) for their images in
\(D\), lift them to \(H_i\in E\), and transport \(T_i\) by
\(H_i\). For a subgroup \(J\) onto \(D\), choose \(j_i\in J\)
above \(h_i\). Since \(H_ij_i^{-1}\in A\), the transported
intersection is exactly
\(j_i(T_i\cap J)j_i^{-1}\); empty intersections remain empty.
The pair of \(j_i\) with the projection of \(r_i\) to
\(Q_\rho\) conjugates the matrices in (E.57.1) literally.
Thus every subgroup homomorphism count is unchanged. Finite subgroup
lattice inversion gives equality of the epimorphism counts as well:
\[
 \#\operatorname{Epi}_{\rm per}(\Lambda,E)
       =\sum_{J\le E}\mu(J,E)\,
          \#\operatorname{Hom}_{\rm per}(\Lambda,J).
\]
Subgroups not onto \(D\) contribute zero. The original map \(q\)
makes the common reference count positive. In the modular case the
conjugacy classes already force the two source orders, so every map
counted is a map from the actual modular source.

Now precompose the source map to \(D^\dagger\) by the hypothesized
automorphism of \(\Lambda/[N,\Lambda]\). The preceding positive
count produces an epimorphism to \(E\) with the prescribed
conjugators modulo \(D\). In the compact space \(\Lambda^3\)
impose the peripheral product equation, generation in each finite
quotient, and agreement with all prescribed conjugator coordinates
modulo \(N\). Every finite collection is covered by one joint finite
quotient \(E\), and is solvable by the count. Their intersection
gives a surjective continuous endomorphism of \(\Lambda\).
It is an automorphism: inverse image under a surjection permutes the
finite set of open subgroups of any given index, so its kernel lies
in every open subgroup and is one. Necessity follows by descending
a full lift; such a lift preserves \(N\), and therefore also
\([N,\Lambda]\).

Finally \(R=R_{\rm nil}(U)\) satisfies \([R,\Lambda]=R\).
In every finite quotient its image is the stable lower central term
of the image of \(U\), by compactness of descending closed terms.
That stable term equals its commutator with the image of \(U\),
hence with the ambient image. Equality in every finite quotient gives
the claimed closed equality. The central criterion therefore applies
directly. The restored independent finite-count program was rerun in
\path{code/check_restored_D50_finite.py}; its finite tests supplement
this proof and do not substitute for the character argument.
\end{proof}

\begin{result}{E.58}{Proved}{LM4 A.18: mixed-sign involutions do not lift with their order}
Here \(F=\widehat F(x,y)\) and \(c(x)=x^{-1},c(y)=y^{-1}\).
The automorphism of \(\prod_pF^{(p)}\) equal to standard conjugation
on \(F^{(3)}\) and the identity on every other factor has a full
peripheral lift, but has no full peripheral involution lift.
Thus ordinary lifting in E.57 cannot preserve order in general.
\end{result}
\begin{proof}
Its multiplier is the genuine unit \(\lambda\) with
\(\lambda_3=-1\) and \(\lambda_p=1\) for \(p\ne3\).
The product action preserves the three peripheral classes and has
order two, as its three-primary abelianization shows. Choose a Galois
automorphism of multiplier \(\lambda\), using the surjectivity of
the cyclotomic character. Translate by its inverse, apply E.57 to the
resulting multiplier-one pronilpotent action, then translate back.
This produces a full ordinary peripheral lift.

We prove directly the only torsion exclusion needed here: a full pure
peripheral involution with \(\lambda_2=1\) is the identity.
A pure peripheral outer class has a unique normalized representative
\(x\mapsto x^\lambda,\ y\mapsto f^{-1}y^\lambda f\), with
\(f\in F'\). Existence follows by adjusting the two abelian
coordinates of a cusp conjugator; uniqueness follows from the
free-factor centralizers. The normalized representatives form a
homomorphic section, so an outer involution gives an actual
automorphism \(\alpha\) with \(\alpha^2=1\).
On \(F^{(2)}\), it is the identity: it is trivial on the full
\(\mathbb Z_2\)-abelianization, and a first nonzero correction in
the Magnus filtration would double on squaring, impossible in a
torsion-free graded piece. The filtration separates the free pro-two
group. Thus \(F\rtimes_\alpha C_2\) maps to
\(F^{(2)}\times C_2\).

For a cusp generator \(c_i\), its normalizer in this semidirect
product is an extension of \(C_2\) by
\(\overline{\langle c_i\rangle}\). Write a lift of the nontrivial
coset as \(a_i\), with \(a_i^2=c_i^{m_i}\).
Then \((\lambda-1)m_i=0\). In the pro-two product quotient,
the free-factor centralizer makes the two-adic coordinate of
\(m_i\) divisible by two. At an odd prime with \(\lambda_p=1\)
one can divide by two, and where \(\lambda_p=-1\) the coordinate
of \(m_i\) is zero. Hence
\((1+\lambda)u_i=-m_i\) is soluble in \(\widehat{\mathbb Z}\),
and \(a_ic_i^{u_i}\) is an involution. Its section subgroup
centralizes \(c_i^{e_2\mathbb Z_2}\), where \(e_2\) is the
dyadic idempotent.

For \(G=F\rtimes_\alpha C_2\), the profinite Brown theorem gives
\[
 \sum_{[S]}\bigl(1+\dim_{\mathbb F_2}H^1(C_F(S),\mathbb F_2)\bigr)
 \le 1+\dim H^1(F,\mathbb F_2)=3. \tag{E.58.1}
\]
Precisely, use \citet[Theorem 2 and \S1.1, equations(1)--(6),
author pp.~3--4]{SymKochloukovaZalesskii2007}. There is no
\(C_2\times C_2\) subgroup, since finite subgroups inject into
\(G/F=C_2\); the open subgroup \(F\) has cohomological
two-dimension one. The theorem's dense restriction embedding is an
isomorphism here because the source has finite dimension. The
normalizer of each order-two subgroup is \(C_F(S)\times S\),
so its K\"unneth formula gives exactly (E.58.1).
Each peripheral section contributes at least two, since the ambient
abelian character detects its centralized dyadic cusp subgroup.
Therefore the three peripheral sections are conjugate. Their common
fixed subgroup \(C\), after these conjugations, contains conjugates
of all three dyadic cusp subgroups. The first two independent ambient
abelian characters show \(\dim H^1(C,\mathbb F_2)\ge2\), and
(E.58.1) makes it exactly two.

We finish with the required rank argument, rather than assuming that
a closed subgroup of rank two must be all of \(F\). If \(C\) were
proper and open, Schreier's formula already gives rank greater than
two. If it has infinite index, the two-adic valuations of
\([F:V]\), for open \(V\supset C\), are unbounded. Otherwise
choose a maximal valuation; all further indices are odd, so
restriction on \(H^1(-,\mathbb F_2)\) is injective by
restriction--corestriction. The unbounded Schreier ranks would embed
in \(H^1(C,\mathbb F_2)=\varinjlim_{V\supset C}H^1(V,\mathbb F_2)\),
a contradiction. This direct-limit identity follows by extending a
finite character from \(C\) to \(CW\) after choosing an open
normal \(W\) whose intersection with \(C\) lies in its kernel.

Choose such a cover of even degree. Each of the three selected
boundaries has odd degree, because its group contains the full
dyadic cusp subgroup. There must be at least a fourth boundary:
otherwise the unique boundary over each base cusp would have the
even total degree. Boundary homology has only the total-sum relation,
so those three classes are independent modulo two. Their restrictions
to \(C\) remain independent, since the selected dyadic elements are
odd multiples of them. This contradicts \(\dim H^1(C,\mathbb F_2)=2\).
Thus \(C=F\). The order-two subgroup defining \(C\) then
centralizes all of \(F\), so \(\alpha\) is inner. Its normalized
representative is consequently the identity.
Every proposed involution lift of the displayed mixed-sign action
would have \(\lambda_2=1\) and \(\lambda_3=-1\), contradicting
this exclusion. No classification of all full peripheral involutions,
nor their conjugacy to complex conjugation, has been imported.
\end{proof}

\begin{result}{E.59}{Proved}{The exact finite A.76 and A.79 detectors, and removal of cusp normalization}
The original stored finite datum in LM4 A.76 excludes a full globally
negative involution retaining its specified finite \(s\)-frame and
literal cusp. The changed stored finite action in A.79 excludes such
a full involution without prescribing either frame: every full positive
peripheral lift fixing the literal cusp has a negative image of exact
order four on the actual finite detector of order
\(2^{2700}72\,7^{13}\). Its old quotient has order
\(2^{972}72\,7^{13}\).

Moreover, for this A.79 action and every prime
\(p\nmid2\cdot72\cdot7\), there is an involution on the actual
quotient \(\mathfrak B/\Phi_p(N_0)\), with both old norm-one
frames retained, that has no full negative peripheral involution lift,
even allowing an unspecified full cusp conjugator.
Here \(N_0\) is the kernel of the old finite quotient and
\(\Phi_p(N_0)=\overline{N_0'N_0^p}\).

These are independent proofs of the negative finite clauses of A.76,
A.79 and the subsequent unnormalized obstruction. They do not by
themselves certify the additional positive ordinary-lifting and
all-iteration clauses of A.79, or the first-central-compatibility
construction of A.81.
\end{result}
\begin{proof}
\emph{The exact source and its coefficient module.}
In this proof the canonical modular involution is
\(c(s)=s\), \(c(r)=sr^{-1}s\), so \(c(\tau)=\tau^{-1}\).
Let \(Q_6=S_3\times A_4\), with
\[
 s=((12),(12)(34)),\qquad r=((123),(123)).
\]
These give a surjection from \(\mathfrak B\). Its kernel is
free of rank thirteen. The actual first seven-primary relation
quotient is
\[
 H=M_7\rtimes Q_6,\qquad M_7=\mathbb F_7^{13}.
\]
The splitting and source vectors are obtained by averaging the
augmented cellular sequence of its Bass--Serre graph: seventy-two
edges, thirty-six order-two vertices, and twenty-four order-three
vertices. Thus this is an actual source quotient, not an arbitrarily
chosen semidirect product.

Put \(k=\mathbb F_8\). The preserved builder
\path{code/lm4_source_verifiers/work_A/check_Q6_forced_frame.py}
constructs the specified nonzero character of \(M_7\), with
\(Q_6\)-orbit of size thirty-six and stabilizer \(\langle s\rangle\).
Its character orbit is disjoint from its two nontrivial Frobenius
translates and is closed under negation. The induced monomial module
\(V=k^{36}\) is consequently simple as an additive binary
\(H\)-module, with endomorphism field exactly \(k\).
The three cyclic fixed spaces have dimensions
\[
 \dim_k V^s=20,\qquad \dim_k V^r=12,\qquad
 \dim_k V^\tau=4,
 \qquad \tau=sr.
\]
The source cocycle equations are precisely
\((1+s)d(s)=0\) and \((1+r+r^2)d(r)=0\).
Their solution space has dimension \(20+24=44\), the coboundary
space has dimension thirty-six, and the complete restriction map is
\[
 H^1(\mathfrak B,V)\xrightarrow{\sim}
 H^1(\langle s\rangle,V)\oplus H^1(\langle\tau\rangle,V),
 \qquad 8=4+4.                                      \tag{E.59.1}
\]
The last cyclic group here is the actual procyclic source subgroup.
Its cohomology is \(V/(\tau-1)V\); it is not replaced by the
cohomology of its finite image.
The independent structural program
\path{code/check_symmetry_source_structure.py} checks the complete
weight orbit, the additive endomorphism field, and all these ranks.

Choose the eight recorded cocycles. Their joint affine image is
all of \(V^8\): a proper residual image would have a nonzero
simple quotient, giving a linear combination that factors through
\(H\), whereas \(H^1(H,V)=0\). This last vanishing, and
\(H^2(H,V)=0\), follow by averaging over the normal odd subgroup
\(M_7\), whose invariants on \(V\) are zero. The old finite
group \(Q\) has normal forms
\[
 (a,m,q),\qquad a\in V^8,\quad
 m\in\operatorname{Sym}^2_k(k^8),\quad q\in H.
\]
Its central cocycle pairs opposite physical weights. When a monomial
permutation reverses an oriented opposite pair, add the corresponding
quadratic carry; the displayed order-two source is repaired by the
recorded parity carry. These formulas are fully preserved in
\path{code/lm4_source_verifiers/work_A/check_Q6_central_cover.py}.
They give \(S^2=R^3=1\). Opposite-weight commutators give all
thirty-six symmetric matrix units, so the actual source image is the
entire displayed group. Its order is
\(2^{3(36\cdot8+36)}72\,7^{13}\), as asserted.

The two witness files are retained byte for byte. For A.79 the hash is
\[
\begin{gathered}\texttt{4135fabf1500006d5ec6bbe5dc0571ea}\\[-2pt]
\texttt{91018abf911e32136cf965729d18d4e9}.
\end{gathered}
\]
The original A.76 hash and both complete hashes occur in
\path{code/check_symmetry_stored_frames.json}.
That independent substitution checks literal cusp preservation,
the negative square, both source conjugacy equations and both
normalizer norms under the original group law. The maps are products
of invertible coefficient maps, the written quadratic carry, inner
automorphisms and cocycle shears. Thus they are automorphisms;
identities on the generating source pair are identities on \(Q\).

\emph{The A.79 detector and forced descent of every candidate.}
Use the four normalized cusp cocycles \(a_i\), for which
\(a_i(s)=0\), and put \(T_0=V^{(3)}\), the cubed monomial
representation, and \(T=T_0^{16}\). Define the actual finite law
\[
\begin{aligned}
 E&=Q\times T,\\
 (g,z)(h,w)&=
 \bigl(gh,z+\rho_3(g)w+
       (a_i(g)^{[2]}\,\rho(g)a_j(h))_{0\le i,j\le3}\bigr).
\end{aligned}
                                                        \tag{E.59.2}
\]
Products and squares inside the last parenthesis are coordinatewise.
The source lifts are \(S=(s,0)\) and
\(R=(r,\mathcal C(a(r),\rho(r)^2a(r)))\); direct multiplication
gives their orders two and three. The image is all of \(E\).
Indeed, at one physical weight, different-row commutators have the
binary span of \((u^2v,v^2u)\), which is all of \(k^2\), and
same-row commutators have the span of \(u^2v+v^2u\), which is
all of \(k\). Their binary ranks, six and three, are independently
checked; alternatively a trace annihilator gives a polynomial with
distinct powers two and four and forces its coefficients to vanish.
These brackets have zero old opposite-weight center. They therefore
attain the sixteen new copies independently. This proves both the
actual source assertion and \(|E|=2^{2700}72\,7^{13}\).

The same odd averaging, followed by the five-term exact sequence for
\(1\to\ker(Q\to H)\to Q\to H\to1\), gives
\[
 H^1(Q,T_0)=\operatorname{Hom}_H(V^8,T_0)=k^8.
\]
The additive intertwiner is explicitly
\(v_\lambda\mapsto v_{-\lambda}^{,4}\).
The eight source classes form the complete source cohomology, so
inflation to \(H^1(\mathfrak B,T_0)\) is an isomorphism.
For \(N_0=\ker(\mathfrak B\to Q)\), its next five-term
sequence therefore makes transgression injective:
\[
 \operatorname{Hom}_Q(N_0,T_0)\hookrightarrow H^2(Q,T_0).
                                                        \tag{E.59.3}
\]
The sixteen kernel maps of (E.59.2) transgress to
\(e_i^{[2]}\smile e_j\). The fixed old positive action fixes
the cusp classes. Equivariance and (E.59.3) force every full lift
of that old action to preserve their joint kernel and to act
identically on \(T\). Canonical conjugation similarly has the
forced cubed monomial action. No kernel-stability condition has
been added to the candidate lifting problem.

For any \(g\in C_Q(s)\), the lift \((g,0)\) centralizes
\(S\): both cup terms with \(s\) vanish. Hence
\(C_E(S)\to C_Q(s)\) is onto, including every head and
residual coordinate. Two peripheral actions of \(E\) over the
fixed old action therefore have \(T\)-conjugate \(S\)-images,
even when neither frame is prescribed. Their difference is a
\(T\)-valued cocycle on \(Q\). Its \(s\)-restriction class
is zero; literal cusp equality makes its cusp restriction zero.
By (E.59.1), every candidate is
\[
 h_E=\operatorname{Inn}(v)h_*,\qquad v\in T^\tau,
                                                        \tag{E.59.4}
\]
where \(h_*\) is the stored action on (E.59.2).
The negative coefficient action fixes \(T^\tau\) pointwise.
Thus, writing \(b_*=h_*c\), every candidate has square
\[
 (\operatorname{Inn}(v)b_*)^2=b_*^2
   =\operatorname{Inn}(t),\qquad
 t_{00,\,\mathrm{weight}\,0}=1.                          \tag{E.59.5}
\]
The restored verifier
\path{code/lm4_source_verifiers/work_A/code/verify_fixed_action_ordinary_square.py}
was rerun. It checks the ranks \(44,40,36\), the square on
both actual source generators, and all 192 binary cusp-inner
directions. The square moves a source generator. Since \(T\)
has exponent two, it is nonidentity of order two; the negative
action therefore has exact order four. This proves the A.79
negative assertion for every candidate, rather than for one choice
of its frames.

\emph{The original A.76 frame.}
Replace (E.59.2) by the six cups \(e_i\smile e_j\),
\(i<j\), with coefficient \(V^{(2)}\). The same averaging,
complete source cohomology and injective transgression force
preservation of this actual six-cup quotient. Retaining the original
old \(s\)-frame and literal cusp again makes every action change
inner by a cusp-fixed coefficient vector. Transporting its frame
with that inner change leaves its normalizer-square residue unchanged.
Every remaining frame change lies in \((V^{(2)})^s\), of
dimension twenty, and changes the residue by the complete norm
\(1+D_s\), of rank eight. Evaluation at physical weight 23
annihilates that entire image. The six original residues, in the
order \(01,02,03,12,13,23\), are
\[
                         (0,0,5,0,5,5)\in k^6.           \tag{E.59.6}
\]
The rerun direct-law verifier checks all 360 binary frame columns
and all 72 cusp-inner columns. This is a linear annihilator of the
complete allowed variation space, not a failure of a search.

If a full negative involution existed with the specified old frame
\(\beta\), then
\(\beta^{-1}b(\beta^{-1})\in C_{\mathfrak B}(s)=\langle s\rangle\),
by the full free-factor theorem
\citep[Theorem B']{SymHerfortRibes1985}. Its old \(S_3\)-image
is one, whereas \(s\) has nonidentity image, so this element is
one. Its six-cup residue would vanish, contradicting (E.59.6).

\emph{An actual odd relation quotient removes the cusp normalization.}
Here is a direct construction, with no appeal to an unspecified
centralizer theorem. Put \(L=N_0\), \(M=L/\Phi_p(L)\),
and let \(b_0=h_0c_Q\) be the checked old involution. Its two
checked norm-one frames allow the permutation modules
\(k_p[Q/\langle s\rangle]\) and
\(k_p[Q/\langle r\rangle]\) to extend to
\(Q\rtimes\langle b_0\rangle\), where \(k_p=\mathbb F_p\).
For instance \(b_0(q\langle s\rangle)=b_0(q)\beta\langle s\rangle\);
the norm-one equation makes its square one. The Bass--Serre sequence
therefore gives a virtual extension of the actual relation module:
\[
 [M]=[1]+[k_pQ]-[k_p[Q/\langle s\rangle]]
                    -[k_p[Q/\langle r\rangle]].         \tag{E.59.7}
\]
Its boundary submodule is the permutation module on
\(Q/\langle t\rangle\), modulo the total sum, with
\(t=q(\tau)\). Give it the negative permutation action of
\(b_0\); this sends \([\tau^{\operatorname{ord}(t)}]\)
to its negative.

For an index-two extension in characteristic prime to its order,
every actual subgroup module whose class is a virtual restriction is
an actual restriction. To prove this, apply semisimplicity to each
orbit of simple subgroup modules. A two-element orbit extends by
induction. For a fixed simple module \(W\), write
\(\operatorname{Ind}W=\bigoplus_i X_i^{n_i}\), with
\(\operatorname{Res}X_i=W^{e_i}\). Reciprocity gives
\(\sum n_ie_i=2\); all possible multiplicities are therefore
one, or the only possible multiplicity is two. Every nonnegative
integer in their restriction lattice is realized. Applying this to
the semisimple complement of the boundary module in (E.59.7) gives
a semilinear involution \(B\) on the actual \(M\), retaining
that negative boundary action.

Odd averaging splits the actual finite extension as \(M\rtimes Q\).
Write its actual cusp as \((a,t)\), and put \(m=\operatorname{ord}(t)\).
The boundary condition is
\(B\bar a=-\bar a\), with
\(\bar a=m^{-1}\sum_{j=0}^{m-1}t^ja\). Thus
\[
 (1-t^{-1})v=-t^{-1}a-Ba
\]
is soluble. Twisting \((B,b_0)\) by \(\operatorname{Inn}(v)\)
inverts the literal cusp. Its square is inner by \(v+Bv\),
which is both \(B\)-fixed and \(t\)-fixed. Subtracting
\((v+Bv)/2\) from \(v\) makes the square one while keeping
the cusp equation. To lift the frames, work in
\((M\rtimes Q)\rtimes\langle b_p\rangle\). The old subgroups
\(\langle s,\beta^{-1}b_0\rangle\cong C_2\times C_2\)
and \(\langle r,J^{-1}b_0\rangle\cong S_3\) have order
prime to \(p\). Averaging in degree two splits their preimages,
and averaging in degree one conjugates each complement so that it
contains the actual source element \(S\), respectively \(R\).
The involution in each complement supplies the required lifted
frame, with its prescribed old image and norm one. The two choices
are independent. This constructs the asserted actual finite
involution for every indicated \(p\).

In this relation quotient,
\[
 C_{M\rtimes Q}(\tau)\longrightarrow Q
       \quad\hbox{has image }\langle t\rangle.           \tag{E.59.8}
\]
Indeed a centralizer element fixes the positive boundary class
\([\tau^m]\). Distinct boundary classes are distinct modulo the
total-sum relation in odd characteristic, and their stabilizer in
\(Q\) is exactly \(\langle t\rangle\). Conversely cusp
powers give that entire image.
In the full group,
\(C_{\mathfrak B}(\tau)=\overline{\langle\tau\rangle}
\cong\widehat{\mathbb Z}\): intersect with the index-six pure
group, where \(\tau^2\) is a free-factor cusp generator, apply
the full free-factor centralizer theorem, and use the order-two
cusp image in \(S_3\).

Suppose a full involution satisfies
\(b(\tau)=u\tau^{-1}u^{-1}\) and induces the constructed
finite action. Equation (E.59.8) gives \(q(u)\in\langle t\rangle\).
Involutivity gives \(b(u)u=\tau^a\). Applying \(b\) to
this equality and also using the cusp equation gives
\(\tau^a=\tau^{-a}\); torsion-freeness gives \(a=0\).
Multiplying \(u\) by a cusp power arranges \(q(u)=1\)
without changing its conjugacy equation or norm. Then
\(\operatorname{Inn}(u^{-1})b\) is a full involution,
inverts the cusp literally, and still induces \(b_0\) on
\(Q\). This contradicts (E.59.5).

\emph{The exact remaining positive comparison.}
If an actual quotient \(\mathfrak B/U\) with
\([U,\mathfrak B]=U\) carries the old action and both norm-one
frames, and maps to \(Q\), then \(U\subset N_j\) for every
\(N_{j+1}=[N_j,\mathfrak B]\). Thus all actual central iterates
inherit the action and frames, and E.57 gives one ordinary full
lift. Each iterate is finite: its central kernel is finitely generated,
and transfer shows its exponent divides six times the preceding
finite index. This deduction is rigorous, but constructing that
particular \(U\) and its action is additional. LM4 does so through
A.56, A.62--65. Their entire integral and older-family compatibility
construction is not certified by the finite calculations above.
Similarly, A.81 additionally constructs an anti-symplectic action on
compact cover homology and solves its actual first central cusp
equation at selected fresh primes. It claims first-central
compatibility, not compatibility at every iterate. Those extra
positive assertions remain separately marked pending in the historical
audit. No failed compactness argument follows from their scope.
\end{proof}
